\emph{Data Flyway}, illustrated in Figure~\ref{fig:transformation-arch}, addresses the five challenges of~\S\ref{sec:introduction} through a single design decision, which is where a pipeline lives. In the prior systems described in~\S\ref{sec:background_related_motivation} it lives in application code, as calls issued at each transfer, so the application holds the ordering of operations and carries intermediate results between them. In Data Flyway it lives in the data management layer, as a directed graph attached to the object or region it applies to. Because the runtime rather than the caller holds that graph, it can record which state each object occupies (\C{2}), see every edge at once rather than one call at a time and so determine when an operation's inputs are complete (\C{3}), choose between alternative routes that a call
site cannot express since a call names one implementation (\C{4}), and traverse the graph when a transfer arrives rather than when a call is written, which places execution in the server and off the application's critical path (\C{1}, \C{5}).

The remainder of this section develops that design. We begin with the architecture and the division of work among its three components (\S\ref{sec:arch}), then the graph representation itself (\S\ref{sec:graph}) and how a client declares a pipeline and attaches it (\S\ref{sec:pipeline}). We then describe how execution overlaps with data movement (\S\ref{sec:overlap}), the cost model that predicts operation latency (\S\ref{sec:model}), and the scheduling algorithm that uses it to choose a device at runtime (\S\ref{sec:sched}).

\subsection{Overview}
\label{sec:arch}

\begin{figure}[t]
    \centering
    \includegraphics[width=\columnwidth, keepaspectratio]{figures_source/images/in_paper/screenshot_2026-09-22.png}
    \caption{Data Flyway architecture within PDC. \blcirc{D} indicates data movement optimizations, \blcirc{P} and \blcirc{C} show computing resources used by clients and PDC respectively, and \blcirc{T} represents the alternative path taken by objects or regions with an attached pipeline.}
    \label{fig:transformation-arch}
\end{figure}

Figure~\ref{fig:transformation-arch} shows Data Flyway as a management layer inside the PDC server, drawn with a light blue background. A client issues region transfers exactly as it would without Data Flyway. Objects carrying an attached pipeline take the alternative path marked \blcirc{T}, along which the layer performs whatever work the transfer requires before data reaches its destination.

Placing the layer in the server rather than in the client library is what allows operations to proceed alongside application computation rather than extending the application's critical path (\C{5}). It also gives the layer access to the full set of devices attached to a node, which is the precondition for choosing among them (\C{4}).

Three components divide the work.

\begin{itemize}
    \item The \textbf{graph optimizer} holds the pipeline attached to each entity and the state that entity currently occupies, and determines which operations an arriving transfer requires and in what order (\C{2}, \C{3}).
    \item The \textbf{scheduler} assigns each operation to a device using predicted cost, accounting for current utilization and for the cost of moving data to that device (\C{4}).
    \item The \textbf{graph executor} runs the operation, records the resulting state, and routes each result to its destination (\C{3}).
\end{itemize}

None of the three is visible to the application, which issues transfers and nothing else (\C{1}).

% Figure~\ref{fig:transformation-arch} shows the PDC components relevant to Data Flyway, which is the in-flight processing management layer shown with a light blue background. Data movement, marked \blcirc{D}, occurs between PDC clients and servers, between the management layer and the computing devices, and between PDC servers and persistent storage such as a parallel file system or cloud storage. HPC memory and storage tiers have grown too numerous for a client to exploit directly, and the same heterogeneity holds across computing devices, where the best place to run a computation depends on dynamic factors including data movement cost and current device utilization. Data Flyway hides this complexity from the client while still using the available hardware. The management layer has three components. The graph optimizer determines a path through the pipeline graph, the scheduler selects a device for each operation on that path, and the graph executor performs the computation on the client's behalf.

\subsection{Creating a Pipeline}
\label{sec:pipeline}

To run operations on a client's behalf, Data Flyway needs a description of
what can be done to the data and under what conditions. For a
transformation this means the states data may occupy and the operations
that move it between them. For an analysis it means the objects an
operation consumes and the objects it produces. Both are captured in one
structure. The graph optimizer operates over a directed graph whose nodes
are data states and named objects, and whose edges are the operations
between them, which lets a client express pipelines of arbitrary shape.

\begin{lstlisting}[language=json, firstnumber=1,
  caption={JSON schema for the transformation graph. The ellipsis \texttt{...} indicates that one or more entries may be present. Fields marked ? are optional.},
  label={fig:transformation-schema},
  float=htbp
]
{
  "lib_path":? "string",
  "name": "string",
  "states": [
    {
      "name": "string"
    },
    ...
  ],
  "transformations": [
    {
      "device": "CPU | GPU",
      "input_state": "states[i].name",
      "output_state": "states[j].name",
      "location": "builtin | external",
      "params": "string",
      "name": "string"
    },
    ...
  ]
}
\end{lstlisting}

\begin{lstlisting}[language=json, firstnumber=1,
  caption={JSON schema for the analysis graph. Fields follow the same conventions as Listing~\ref{fig:transformation-schema}. The \texttt{input\_objects} and \texttt{output\_objects} fields are arrays rather than single names, which is what expresses fan-in and fan-out.},
  label={fig:analysis-schema},
  float=htbp
]
{
  "lib_path":? "string",
  "name": "string",
  "objects": [
    {
      "name": "string"
    },
    ...
  ],
  "analyses": [
    {
      "device": "CPU | GPU",
      "input_objects": ["objects[i].name", ...],
      "output_objects": ["objects[j].name", ...],
      "location": "builtin | external",
      "trigger": "eager | lazy",
      "persist": "true | false",
      "params": "string",
      "name": "string"
    },
    ...
  ]
}
\end{lstlisting}

A pipeline may contain several paths between the same pair of states, for
example CPU and GPU implementations of one transformation, which makes an
imperative API awkward to build around. Data Flyway therefore describes the
graph in a JavaScript Object Notation (JSON) file. JSON is human readable,
widely supported, and represents graph structures naturally as arrays of
objects, so it needs no custom parser and fits existing HPC workflow
tooling.

Listings~\ref{fig:transformation-schema} and~\ref{fig:analysis-schema} give
the two schemas. When a client passes a file path to PDC, Data Flyway
validates the file for well-formedness and semantic correctness. The
optional \texttt{lib\_path} field names a shared library holding
user-defined implementations, and \texttt{name} identifies the pipeline. In
the transformation schema, the \texttt{states} array declares the
intermediate states data may occupy, each with a unique name, and the
\texttt{transformations} array declares the operations that convert between
them. Each entry gives its execution device, its input and output states,
and whether it is built in or supplied through a shared library.
Configuration parameters pass through \texttt{params}, and \texttt{name}
identifies the function to run.

The analysis schema is analogous but operates over \texttt{objects} rather
than \texttt{states}. The \texttt{objects} array declares the named objects
an operation may consume or produce, and the \texttt{analyses} array
declares the operations that derive objects from other objects. Because one
operation may take several inputs and produce several outputs, its
\texttt{input\_objects} and \texttt{output\_objects} fields are arrays,
which is how a pipeline expresses fan-in and fan-out directly in the
description.

Evaluation timing follows from two independent fields rather than from any
special-cased execution path. Setting \texttt{trigger} to \texttt{eager}
makes the executor compute an operation as soon as its last required input
is written. Setting it to \texttt{lazy} defers computation until a client
reads an output object, recomputing from the current inputs on every such
read. The \texttt{persist} field independently controls whether a computed
result is stored as an ordinary object. Eager with \texttt{persist: true},
the default, gives eager evaluation. Lazy with \texttt{persist: false}
gives lazy evaluation. Eager with \texttt{persist: false} gives the
transient intermediates that chain analysis stages, such as
\texttt{curl\_x}, \texttt{curl\_y} and \texttt{curl\_z}
in~\S\ref{sec:curl-vorticity-eval}, which downstream stages consume without
their ever reaching storage. Because both fields are set per operation,
eager, lazy and transient stages coexist in one pipeline.

Once a graph is attached to a PDC entity, the server-side management layer
handles every subsequent transfer. PDC object metadata records which
entities carry an attached graph, what state each currently holds or which
objects have already been derived from it, and timing from the operations
that produced that state. The graph optimizer and the scheduler read this
metadata to decide how to reach a requested state, what to derive, and
where and when to move data and execute each operation. Subsequent calls to
open an entity and transfer data need no knowledge of the graph. Data
Flyway determines from the metadata whether a transformation or a
derivation is required, and the executor delivers the data in the expected
state without client involvement.

A client's workflow is therefore as follows.

\begin{enumerate}
    \item Create a JSON file describing a pipeline.
    \item Register the file with the data management system at process initialization.
    \item Attach pipelines to the relevant objects or regions.
    \item Perform core computation and region transfers.
    \item Free pipeline resources.
\end{enumerate}

Only steps 4 belongs to the application's steady state. After attachment, a
developer issues ordinary region transfers and nothing else. The API also
provides debugging utilities, including a check that a valid path exists
between two states, or that every input an analysis operation requires is
reachable in the graph.

\subsection{Predictive Model for GPU Kernel Latency}
\label{sec:model}

We predict total GPU kernel latency ($\texttt{total\_ms}$) using polynomial
regression over NVML hardware counters together with \emph{lag features},
which capture autocorrelation across successive invocations. The lag
features are the host-to-device and device-to-host latency of the
immediately preceding invocation ($\texttt{prev\_h2d\_ms}$,
$\texttt{prev\_d2h\_ms}$). This extends the model
of~\cite{ravi2023runway}, which achieves $R^2 = 0.50$ on A100s from payload
size and GPU utilization alone. As Figure~\ref{fig:runway} shows, that
baseline yields $R^2 = 0.38$ on our platform, also using A100s, with
predictions collapsing into a narrow band that fails to span the observed
latency range. Forward-backward BIC stepwise regression selects six base
predictors ($\texttt{data\_size\_mb}$, $\texttt{prev\_h2d\_ms}$,
$\texttt{prev\_d2h\_ms}$, $\texttt{nvml\_gpu\_util}$,
$\texttt{nvml\_pcie\_tx\_kbps}$, $\texttt{nvml\_pcie\_rx\_kbps}$), which are
expanded to a degree-2 polynomial, and iterative term-level pruning at $p 
10^{-4}$ reduces that expansion to 15 retained monomials. We split the $n =
2{,}320$ observations 80/20, using the training portion for both predictor
selection and fitting and holding out the remainder for evaluation. As
Figure~\ref{fig:bestr2} shows, the resulting model achieves $R^2 = 0.97$ on
the held-out set, with predictions distributed tightly around the diagonal
across the full 0--160\,ms range.

\begin{figure}[t]
  \centering
  \includegraphics[width=0.4\columnwidth]{figures_source/images/in_paper/residual_legend.pdf}

  \begin{subfigure}[b]{0.48\columnwidth}
    \includegraphics[width=\linewidth]{figures_source/images/in_paper/runway_predicted_vs_actual.pdf}
    \caption{Baseline ($R^2=0.38$).}
    \label{fig:runway}
  \end{subfigure}
  \hfill
  \begin{subfigure}[b]{0.48\columnwidth}
    \includegraphics[width=\linewidth]{figures_source/images/in_paper/best_r2_predicted_vs_actual.pdf}
    \caption{Our model ($R^2=0.97$).}
    \label{fig:bestr2}
  \end{subfigure}
  \caption{Predicted against actual total kernel latency. The baseline, which uses only payload size and GPU utilization, fails to capture latency variance across workloads. Our 15-term degree-2 polynomial augmented with lag features tracks the diagonal across the full 0--160\,ms range. Both panels report $R^2$ on the held-out set.}
  \label{fig:models}
\end{figure}

\subsection{Dynamic Scheduling Algorithm}
\label{sec:sched}

When a client transfers data, Data Flyway must decide at runtime which
device executes each operation. The decision is difficult because device
utilization fluctuates, operation cost varies with data size and type, and
moving data between devices carries its own overhead. A static assignment
cannot account for contention when devices are saturated, nor for cases
where transfer cost outweighs the benefit of the faster device. The
algorithm described here is the core of the scheduler, which works with the
graph optimizer and the executor to select both a path and a device
assignment and then run the operation.

\begin{algorithm}[t]
\caption{Dynamic scheduling algorithm.}
\label{fig:scheduling_algorithm}
\begin{algorithmic}[1]
\STATE $edge \gets \text{get\_next\_best\_edge}($
\STATE \hspace{2em} $dg,\ cur\_state,\ desired\_state)$
\WHILE{$edge \neq \text{null}$}
    \STATE $data\_size \gets \text{bytes}(input\_region)$
    \STATE $move_{CPU} \gets \text{movement}(device,\ \text{CPU})$
    \STATE $move_{GPU} \gets \text{movement}(device,\ \text{GPU})$
    \STATE $h2d \gets (move_{GPU} > 0)\ ?\ prev\_h2d : 0$
    \STATE $d2h \gets (move_{CPU} > 0)\ ?\ prev\_d2h : 0$
    \STATE $cost_{CPU} \gets \hat{f}_{CPU}(data\_size,\ CPU_{util})$
    \STATE $cost_{GPU} \gets \hat{f}_{GPU}(data\_size,\ GPU_{util},\ h2d,\ d2h,$
    \STATE \hspace{2em} $nvml\_pcie\_tx\_kbps,\ nvml\_pcie\_rx\_kbps)$
    \IF{$cost_{CPU} \leq cost_{GPU}$}
        \STATE $device \gets \text{CPU}$
    \ELSE
        \STATE $device \gets \text{GPU}$
    \ENDIF
    \STATE $total\_time \gets \text{run}(edge,\ device)$
    \IF{$device = \text{GPU}$}
        \STATE $\text{update\_lag\_features}_{GPU}(h2d,\ d2h)$
    \ENDIF
    \STATE $edge \gets \text{get\_next\_best\_edge}($
    \STATE \hspace{2em} $dg,\ cur\_state,\ desired\_state)$
\ENDWHILE
\end{algorithmic}
\end{algorithm}

Algorithm~\ref{fig:scheduling_algorithm} schedules operations along the
lowest-predicted-cost path between a region's current state and its desired
state. It begins by calling \texttt{get\_next\_best\_edge}, which computes
the shortest path through the graph and returns the first edge to execute.
Each edge carries as its weight the lower of the CPU and GPU runtimes
predicted by the models below, including the cost of moving data to that
device, so path selection accounts for both computation and data placement
across all remaining steps. A null return means the current state already
equals the desired state and traversal is complete.

For each edge the scheduler first computes the movement cost for each
device, which is the cost of transferring data to the CPU or GPU if it is
not already resident there. Data already on the target device costs
nothing to move, and the corresponding lag feature is set to zero since no
transfer will occur. The scheduler then predicts runtime with
device-specific polynomial models. The CPU model $\hat{f}_{CPU}$ takes data
size and current CPU utilization, while the GPU model $\hat{f}_{GPU}$ also
incorporates the lag features from the previous execution on that GPU to
account for transfer overhead. The operation is dispatched to whichever
device carries the lower total cost.

After execution only the GPU lag features are updated, since only the GPU
model depends on them. The scheduler then calls
\texttt{get\_next\_best\_edge} again, recomputing the best next step from
the updated graph state, which lets it adapt to conditions that change
between operations.